\documentclass[12pt]{article}
% Préambule commun à tous les documents extraits de « La Revue CAPES »
\usepackage[T1]{fontenc}
\usepackage[utf8]{inputenc}
\usepackage{lmodern}
\usepackage{amsmath,amssymb,amsthm,amsfonts,mathrsfs}
\usepackage{setspace}
\usepackage{listings}
\usepackage{pgfplots}
\pgfplotsset{compat=1.18}
\usepackage{longtable}
\usepackage{adjustbox}
\usepackage{lettrine}
\usepackage{tkz-tab}
\usepackage{changepage}
\usepackage{pdflscape}
\usepackage{graphicx}
\usepackage{tikz}
\usepackage{xcolor}
\usepackage[a4paper,top=2cm,bottom=2.5cm,left=2.5cm,right=2cm]{geometry}
\usepackage{fancyhdr}
\usepackage{draftwatermark}
\SetWatermarkText{La RevueCapes}
\SetWatermarkLightness{0.9}
\SetWatermarkScale{2}
\usepackage{hyperref}
\hypersetup{colorlinks=true,linkcolor=blue!50!black,urlcolor=blue!50!black}

\definecolor{revue}{RGB}{20,60,120}

\pagestyle{fancy}
\fancyhf{}
\renewcommand{\headrulewidth}{0.4pt}
\setlength{\headheight}{15pt}

% \entete{Matière}{Titre}{Type de document}
\newcommand{\entete}[3]{%
  \fancyhead[L]{\small\textsc{La Revue CAPES} -- #1}%
  \fancyhead[R]{\small #2}%
  \fancyfoot[C]{\thepage}%
  \thispagestyle{plain}%
  \begin{center}
    {\color{revue}\rule{\textwidth}{1.2pt}}\\[4pt]
    {\small\textsc{Concours directs CAP-CEG / CAPES -- Burkina Faso}}\\[6pt]
    {\LARGE\bfseries\color{revue} #2}\\[6pt]
    {\large\itshape #1 -- #3}\\[2pt]
    {\color{revue}\rule{\textwidth}{1.2pt}}
  \end{center}
  \vspace{0.5cm}%
}

% Encadré pour les remarques ajoutées dans les corrigés
\newcommand{\remarque}[1]{\par\medskip\noindent\fcolorbox{revue}{revue!6}{\parbox{\dimexpr\linewidth-2\fboxsep-2\fboxrule}{\textbf{\color{revue}Remarque.} #1}}\par\medskip}


\begin{document}
\entete{Mathématiques}{Ancien sujet CAPES -- Session 2020}{Énoncé}
\begin{spacing}{1.5}

\textbf{Exercice 1 : (6pts)}

On donne la matrice $ A= \begin{pmatrix}
1 & 0 & -1\\
0 & 1 & 0\\
-1 & 1 & 1
\end{pmatrix}$                                                                   

1. Justifier que la matrice $A$ est inversible.

2. Déterminer le polynôme caractéristique de $A$.

3. Soit $\alpha, \beta$ et $\gamma$ les valeurs propres de la matrice $A$. Déterminer $\alpha, \beta$ et $\gamma$ sachant qu'elles sont classées dans l'ordre croissant.

4. Déterminer les vecteurs propres $u$, $v$ et $w$ associés aux valeurs propres $\alpha, \beta$ et $\gamma$.

5. Soit $D=\begin{pmatrix}
\alpha & 0 & 0 \\ 
0 & \beta & 0 \\ 
0 & 0 & \gamma
\end{pmatrix} $. Donner la matrice de passage $P$ de $A$ à $D$.

\textbf{Exercice 2 : (8pts)}

1. On considère la fonction $f$ définie sur $\mathbb{R}$ par $f(x)=e^{-x}$.

a. Montrer que $f$ est dérivable sur $\mathbb{R}$.

b. Déterminer $f^\prime(x)$, $f^{\prime\prime}(x)$ et $f^{(3)}(x)$; $\forall {x}\in \mathbb{R}$.

c. Démontrer par récurrence que : $f^{(n)}(x)=(-1)^ne^{-x}$



d. Donner le développement limité d'ordre $n$ de $f$ au voisinage de $0$.

2. On considère la courbe paramétrée : $(\Gamma) :\begin{cases}
x(t)=\cos 3t\\
y(t)=\sin 2t
\end{cases} \forall t\in \mathbb{R}$

a. Montrer que l'étude de la courbe peut être réduite à l'intervalle $[0,\pi]$.

b. Dresser le tableau de variation commun de $x$ et $y$ sur $[0,\pi]$.

\textbf{Exercice 3 : (3pts)}

Soient $n\geq 1$, un entier naturel et $\mathbb{R}_n[X]$ l'espace vectoriel des polynômes réels de degré inférieur ou égal à $n$. $\forall {P,Q\in \mathbb{R}_n[x]}$, on pose :

\[ \beta(P,Q)=\int_{0}^{1}xP(x)Q^{\prime}(x). dx\]

1. Montrer que $\beta$ est une forme bilinéaire.

2. Exprimer $\beta(x,1)$ et $\beta(1,x)$.

3. $\beta$ est-elle symétrique? Antisymétrique?

\textbf{Exercice 4 :(3pts)}

1. Montrer à l'aide du théorème de Fubini que : 

$I=\int_{0}^{+\infty}e^{-\frac{x^2}{2}}dx=\sqrt{\frac{\pi}{2}}$ (indication : introduire $I^2$).

2. On pose $\forall {n\in \mathbb{N}}, \forall {x\in \mathbb{R}}$ $\\
f_n(x)=(1+\frac{x^2}{n})-\frac{n+1}{2}$ et $a_n=\int_{-\infty}^{+\infty}f_n(x)dx.$

a. Montrer que, $\forall {u\in \mathbb{R}}$ et $\forall {p\in \mathbb{N}}, (1+u)^p\geq 1+pu$.

b. Déduisez que $\forall {x\in \mathbb{R}}, f_n(x) \leq \frac{2}{1+x^2}$.

c. Calculer $lim_{n\to +\infty}a_n$.

\end{spacing}
\end{document}
